\section{Analysis}
\label{sec:analysis}

\subsection{When Does Task-Grounded Scoring Excel?}

Our method's performance correlates strongly with the \emph{information density} of item textual representations. We quantify this through average title length and description richness:

\begin{itemize}
    \item \textbf{High information density} (Books: 49-char titles, 782-char descriptions): NDCG@10 improvement of +21.5\% over best baseline.
    \item \textbf{Medium information density} (Electronics: 15-char titles but detailed technical descriptions): +52.5\% improvement.
    \item \textbf{Low information density} (Movies: 21-char titles, 387-char descriptions): Competitive but not SOTA.
\end{itemize}

This pattern is expected: task-grounded scoring relies on the LLM's ability to reason about user-item relevance from text. When text carries rich semantic information, the LLM can make fine-grained distinctions. When text is sparse (e.g., movie titles like ``Shock Waves''), collaborative signals captured by embedding methods become more valuable.

\subsection{Score Distribution Analysis}

We analyze the distribution of relevance probability scores across domains:

\begin{itemize}
    \item \textbf{Books} (SOTA domain): Positive items receive median score significantly higher than negatives, with clear separation. The model confidently identifies relevant items.
    \item \textbf{Movies} (weakest domain): Score distributions for positive and negative items overlap substantially. The model assigns similar probabilities to many candidates, reducing discriminative power.
\end{itemize}

This confirms that the performance gap is not due to calibration failure but to \emph{insufficient discriminative signal} in the textual features available to the model.

\subsection{Comparison with Post-Hoc Calibration}

A natural question is whether existing baselines could achieve similar calibration through post-hoc methods. We argue this is insufficient:

\begin{enumerate}
    \item \textbf{Temperature scaling} is a monotonic transformation that preserves rank order — it improves ECE but cannot improve HR@$k$ or NDCG@$k$.
    \item \textbf{Isotonic regression} can improve calibration but requires a held-out calibration set and does not improve the underlying discriminative quality.
    \item \textbf{Our approach} improves both calibration \emph{and} ranking simultaneously because it changes the scoring mechanism itself, not just the score distribution.
\end{enumerate}

\subsection{Computational Cost}

Our method requires one LLM forward pass per user-candidate pair, identical to pointwise scoring baselines. With vLLM batch inference (batch size 512, prefix caching enabled), we process approximately 38 samples/second on a single RTX 4090, completing 10,000 users $\times$ 101 candidates in approximately 7.4 hours. This is comparable to other LLM-based methods that require per-item scoring.

\subsection{Limitations}

\begin{enumerate}
    \item \textbf{Text dependency:} Performance degrades on domains with uninformative item descriptions. Hybrid approaches combining our scoring with collaborative signals could address this.
    \item \textbf{Single backbone:} We evaluate only on Qwen3-8B. Generalization to other LLM families (LLaMA, GPT) requires further study.
    \item \textbf{Candidate set size:} We evaluate with 101 candidates. Scaling to full-catalog recommendation requires efficient candidate generation as a first stage.
    \item \textbf{Single run:} Results are from single inference runs without multiple seeds. Statistical significance tests with bootstrap confidence intervals are needed for publication.
\end{enumerate}
